\scanpage{163}
Si $\widetilde\xi$ \deleted{est} réelle, on a une structure de Hodge à involution.

Comment se relient ces structures de Hodge ??

\diagram{\begin{tikzcd}[column sep=huge,row sep=large]
\substack{\widetilde{\mathcal M}_0\\\text{motifs}\\\text{abéliens sur}\\\widetilde K\ \text{provenant}\\\text{de }K} \arrow[r,"{\text{fibre en }\widetilde\xi}"] & \deleted{\text{Vectoriels}}\ G^0\text{-Modules} \arrow[r,"\Phi_{\widetilde\xi}"] & \substack{\text{Réseaux de}\\\text{Hodge}} \\
{} \arrow[r,"{\text{fibre en }\widetilde\xi'}"'] & {} \arrow[r,no head,"\Phi"'] & {}
\end{tikzcd}}

\deleted{\unsure{NB} La donnée de $\xi$}

\par\noindent
\deleted{\mbox{\begin{minipage}[c]{\linewidth}
\uline{NB} Le seul \deleted{bon} caractère \unsure{invariant|rationnel} de $\uuline{K}^*$ est $N_{K^*/\mathbf Q}$ ; ainsi \deleted{les}
\diagram{\begin{tikzcd}[column sep=large]
\uuline{K}^*\text{-modules} \arrow[r,"\sim"] & \substack{\text{\unsure{systèmes} de}\\\text{vectoriels de dim finie sur }K}
\end{tikzcd}}
\deleted{modules} Si $E^0\ne0$, l'opération est triviale pour $E^0$ ; le \unsure{sous-espace|seul poids} est nul. Donc les vectoriels sur lesquels $\uuline{K}^*$ opère sont les \deleted{vectoriels} espaces \inserted{vectoriels} gradués par $\mathbf Z$.
\end{minipage}}}
\par

\par\noindent
\deleted{\mbox{\begin{minipage}[c]{\linewidth}
$V$ un $G^0$-module

$V\otimes_{\mathbf Q}\mathbf C$ muni des \deleted{\unsure{une}} structure de Hodge --- $\pi_\xi^V$

($\xi$ place de $\widetilde K$)

et d'un sous-$\widetilde K$-espace vectoriel $H_\xi^V$
\end{minipage}}}
\par
\margin{162}
